\textbf{Phần 1. Trắc nghiệm lựa chọn}
\setcounterpageref{ex}{0}

\begin{ex} %[3D2V78]
Tính giá trị của phép tính: $16\,250 \times 4$.
\choice
{$64\,000$}
{\True $65\,000$}
{$66\,000$}
{$65\,500$}
\loigiai{
$16\,250 \times 4 = 65\,000$.
\textbf{Đáp án của câu hỏi là B}
}
\end{ex}

\begin{ex} %[3D2V78]
Tính giá trị của phép chia: $84\,648 : 4$.
\choice
{$21\,160$}
{$21\,262$}
{$21\,062$}
{\True $21\,162$}
\loigiai{
$84\,648 : 4 = 21\,162$.
\textbf{Đáp án của câu hỏi là D}
}
\end{ex}

\begin{ex} %[3D2V78]
Tìm số tự nhiên $y$, biết: $y : 5 = 18\,000$.
\choice
{\True $90\,000$}
{$85\,000$}
{$95\,000$}
{$80\,000$}
\loigiai{
$y = 18\,000 \times 5 = 90\,000$.
\textbf{Đáp án của câu hỏi là A}
}
\end{ex}

\begin{ex} %[3D2V78]
Tính giá trị của biểu thức: $A = 14\,000 \times 3 + 36\,000 : 6$.
\choice
{$45\,000$}
{$50\,000$}
{\True $48\,000$}
{$46\,000$}
\loigiai{
$14\,000 \times 3 = 42\,000$; $36\,000 : 6 = 6000$; $42\,000 + 6000 = 48\,000$.
\textbf{Đáp án của câu hỏi là C}
}
\end{ex}

\begin{ex} %[3D2V78]
Một hình vuông có chu vi $68\,400\text{ m}$. Độ dài mỗi cạnh hình vuông là
\choice
{$17\,000\text{ m}$}
{\True $17\,100\text{ m}$}
{$17\,200\text{ m}$}
{$16\,100\text{ m}$}
\loigiai{
$68\,400 : 4 = 17\,100\text{ (m)}$.
\textbf{Đáp án của câu hỏi là B}
}
\end{ex}

\begin{ex} %[3D2V78]
Tính nhanh giá trị: $20\,000 \times 3 + 20\,000 \times 2$.
\choice
{$80\,000$}
{$90\,000$}
{$120\,000$}
{\True $100\,000$}
\loigiai{
$20\,000 \times (3 + 2) = 20\,000 \times 5 = 100\,000$.
\textbf{Đáp án của câu hỏi là D}
}
\end{ex}

\begin{ex} %[3D2V78]
Trong phép chia cho 9 có số dư lớn nhất là 8 và thương là $10\,500$. Số bị chia là
\choice
{\True $94\,508$}
{$94\,500$}
{$94\,507$}
{$94\,509$}
\loigiai{
Số bị chia: $9 \times 10\,500 + 8 = 94\,500 + 8 = 94\,508$.
\textbf{Đáp án của câu hỏi là A}
}
\end{ex}

\begin{ex} %[3D2V78]
Một kho gạo có $96\,000\text{ kg}$ gạo. Đã bán đi $\dfrac{1}{4}$ số gạo. Số gạo còn lại là
\choice
{$24\,000\text{ kg}$}
{$70\,000\text{ kg}$}
{\True $72\,000\text{ kg}$}
{$75\,000\text{ kg}$}
\loigiai{
Đã bán: $96\,000 : 4 = 24\,000\text{ kg}$. Còn lại: $96\,000 - 24\,000 = 72\,000\text{ (kg)}$.
\textbf{Đáp án của câu hỏi là C}
}
\end{ex}

\begin{ex} %[3D2V78]
Tìm số tự nhiên $x$, biết: $x \times 4 + 15\,000 = 75\,000$.
\choice
{$12\,000$}
{\True $15\,000$}
{$16\,000$}
{$18\,000$}
\loigiai{
$x \times 4 = 75\,000 - 15\,000 = 60\,000 \Rightarrow x = 60\,000 : 4 = 15\,000$.
\textbf{Đáp án của câu hỏi là B}
}
\end{ex}

\begin{ex} %[3D2V78]
Một mảnh đất hình chữ nhật có chiều rộng $12\,000\text{ m}$, chiều dài gấp đôi chiều rộng. Chu vi mảnh đất đó là
\choice
{$48\,000\text{ m}$}
{$60\,000\text{ m}$}
{$84\,000\text{ m}$}
{\True $72\,000\text{ m}$}
\loigiai{
Chiều dài: $24\,000\text{ m}$. Chu vi: $(24\,000 + 12\,000) \times 2 = 72\,000\text{ (m)}$.
\textbf{Đáp án của câu hỏi là D}
}
\end{ex}

\begin{ex} %[3D2V78]
Thực hiện phép chia: $100\,000 : 4$. Kết quả là
\choice
{\True $25\,000$}
{$20\,000$}
{$30\,000$}
{$22\,500$}
\loigiai{
$100\,000 : 4 = 25\,000$.
\textbf{Đáp án của câu hỏi là A}
}
\end{ex}

\begin{ex} %[3D2V78]
Một xưởng dệt mỗi ngày dệt được $14\,500\text{ m}$ vải. Hỏi trong 6 ngày xưởng dệt được bao nhiêu mét vải?
\choice
{$85\,000\text{ m}$}
{$86\,000\text{ m}$}
{\True $87\,000\text{ m}$}
{$88\,000\text{ m}$}
\loigiai{
$14\,500 \times 6 = 87\,000\text{ (m)}$.
\textbf{Đáp án của câu hỏi là C}
}
\end{ex}

\textbf{Phần 2. Trắc nghiệm đúng sai}
\setcounterpageref{ex}{0}

\begin{ex} %[3D2V78]
Xét tính đúng sai của các khẳng định về phép nhân chia:
\choiceTF
{\True Nhân một số với 1 luôn bằng chính nó.}
{Một số chia cho 1 thì bằng 0.}
{\True Trong biểu thức chỉ có nhân và chia, thực hiện từ trái sang phải.}
{$25\,000 \times 4 = 10\,000$.}
\loigiai{
\begin{itemize}
    \item Ý a) Đúng: Tính chất của 1.
    \item Ý b) Sai: Sai vì chia cho 1 bằng chính nó.
    \item Ý c) Đúng: Thứ tự tính.
    \item Ý d) Sai: Sai vì $25\,000 \times 4 = 100\,000$.
\end{itemize}
\begin{itemize}
    \item a) \textbf{Đúng}.
    \item b) \textbf{Sai}.
    \item c) \textbf{Đúng}.
    \item d) \textbf{Sai}.
\end{itemize}
}
\end{ex}

\begin{ex} %[3D2V78]
Một nhà máy có 4 dây chuyền sản xuất, mỗi dây chuyền làm được 18 000 sản phẩm mỗi tháng:
\choiceTF
{\True Cả 4 dây chuyền làm được 72 000 sản phẩm mỗi tháng.}
{\True Hai dây chuyền làm được 36 000 sản phẩm mỗi tháng.}
{Cả nhà máy làm trong 2 tháng được 150 000 sản phẩm.}
{\True Nếu đóng 6 sản phẩm vào 1 hộp thì mỗi tháng đóng được 12 000 hộp.}
\loigiai{
\begin{itemize}
    \item Ý a) Đúng: $18\,000 \times 4 = 72\,000$ sản phẩm.
    \item Ý b) Đúng: $18\,000 \times 2 = 36\,000$ sản phẩm.
    \item Ý c) Sai: Sai vì $72\,000 \times 2 = 144\,000 < 150\,000$.
    \item Ý d) Đúng: $72\,000 : 6 = 12\,000$ hộp.
\end{itemize}
\begin{itemize}
    \item a) \textbf{Đúng}.
    \item b) \textbf{Đúng}.
    \item c) \textbf{Sai}.
    \item d) \textbf{Đúng}.
\end{itemize}
}
\end{ex}

\textbf{Phần 3. Trả lời ngắn}
\setcounterpageref{ex}{0}

\begin{ex} %[3D2V78]
Tính: $22\,500 \times 4$.
\shortans{90000}
\loigiai{
\textbf{Đáp số:} $90000$\\
$22\,500 \times 4 = 90\,000$.
}
\end{ex}

\begin{ex} %[3D2V78]
Tính: $75\,500 : 5$.
\shortans{15100}
\loigiai{
\textbf{Đáp số:} $15100$\\
$75\,500 : 5 = 15\,100$.
}
\end{ex}

\begin{ex} %[3D2V78]
Tìm $x$, biết: $x \times 6 = 96\,000$.
\shortans{16000}
\loigiai{
\textbf{Đáp số:} $16000$\\
$x = 96\,000 : 6 = 16\,000$.
}
\end{ex}

\begin{ex} %[3D2V78]
Tính giá trị biểu thức: $15\,000 \times 5 - 25\,000$.
\shortans{50000}
\loigiai{
\textbf{Đáp số:} $50000$\\
$75\,000 - 25\,000 = 50\,000$.
}
\end{ex}

\begin{ex} %[3D2V78]
Một hình vuông có cạnh dài $18\,000\text{ m}$. Tính chu vi hình vuông theo mét.
\shortans{72000}
\loigiai{
\textbf{Đáp số:} $72000$\\
$18\,000 \times 4 = 72\,000\text{ (m)}$.
}
\end{ex}

\begin{ex} %[3D2V78]
Có 5 xe chở hàng, mỗi xe chở $14\,000\text{ kg}$ gạo. Hỏi cả đoàn xe chở bao nhiêu ki-lô-gam gạo?
\shortans{70000}
\loigiai{
\textbf{Đáp số:} $70000$\\
$14\,000 \times 5 = 70\,000\text{ (kg)}$.
}
\end{ex}
